% GENERATED FILE. Edit canonical lessons, then run npm run generate:books.
% canonical-source-hash: fnv1a64:6ae583ddcfe5e68a
% canonical-lessons: TA-C243-kotumai, TA-C243-urulaikkilanku, TA-C243-kattarikkay, TA-C243-vacanai, TA-C243-vellarikkay

\chapter{Wheat, Potato, Aubergine, Fragrance, Cucumber}
\label{ch:ta-wheat-potato-aubergine-fragrance-cucumber}
\hlchaptermodality{243}

\begin{chapteropening}
\textbf{By the end of this chapter:} \emph{I can say wheat, a potato, an aubergine, a fragrance and a cucumber.}

Complete the last lesson of chapter 243: I can say wheat, a potato, an aubergine, a fragrance and a cucumber.
\end{chapteropening}

\section[\texorpdfstring{\ta{கோதுமை} (kōtumai)}{கோதுமை (kōtumai)}]{\ta{கோதுமை} (kōtumai) — wheat}

\label{lesson:TA-C243-kotumai}



\begin{quote}
Before the new one: say the Tamil for a button, then the Tamil for a blouse (worn with a sari).
\end{quote}

\subsection*{You'll want to know: \ta{கோதுமை}}
\textbf{\ta{கோதுமை}} — \emph{kōtumai} — \textquotedblleft{}wheat\textquotedblright{}.

\begin{grammarlens}[title={What you've built}]
One more word for this chapter.
\end{grammarlens}

\subsection*{Your turn}
\begin{itemize}
\raggedright
  \item \emph{Say it:} \emph{kōtumai}
  \item \emph{Say it:} \emph{kōtumai}, once more
  \item \emph{Say it:} say \emph{kōtumai}
  \item \emph{Recall:} say the Tamil for a button, then the Tamil for a blouse (worn with a sari), then say \emph{kōtumai} again
\end{itemize}

\begin{tcolorbox}[breakable,colback=teal!4,colframe=teal!35!black,title={Before you move on}]
What does \ta{கோதுமை} mean? (\textquotedblleft{}Wheat\textquotedblright{}.) Say it once more.
\end{tcolorbox}
\section[\texorpdfstring{\ta{உருளைக்கிழங்கு} (uruḷaikkiḻaṅku)}{உருளைக்கிழங்கு (uruḷaikkiḻaṅku)}]{\ta{உருளைக்கிழங்கு} (uruḷaikkiḻaṅku) — a potato}

\label{lesson:TA-C243-urulaikkilanku}



\begin{quote}
Before the new one: say the Tamil for a blouse (worn with a sari), then the Tamil for wheat.
\end{quote}

\subsection*{You'll want to know: \ta{உருளைக்கிழங்கு}}
\textbf{\ta{உருளைக்கிழங்கு}} — \emph{uruḷaikkiḻaṅku} — \textquotedblleft{}a potato\textquotedblright{}.

\begin{grammarlens}[title={What you've built}]
One more word for this chapter.
\end{grammarlens}

\subsection*{Your turn}
\begin{itemize}
\raggedright
  \item \emph{Say it:} \emph{uruḷaikkiḻaṅku}
  \item \emph{Say it:} \emph{uruḷaikkiḻaṅku}, once more
  \item \emph{Say it:} say \emph{uruḷaikkiḻaṅku}
  \item \emph{Recall:} say the Tamil for a blouse (worn with a sari), then the Tamil for wheat, then say \emph{uruḷaikkiḻaṅku} again
\end{itemize}

\begin{tcolorbox}[breakable,colback=teal!4,colframe=teal!35!black,title={Before you move on}]
What does \ta{உருளைக்கிழங்கு} mean? (\textquotedblleft{}A potato\textquotedblright{}.) Say it once more.
\end{tcolorbox}
\section[\texorpdfstring{\ta{கத்தரிக்காய்} (kattarikkāy)}{கத்தரிக்காய் (kattarikkāy)}]{\ta{கத்தரிக்காய்} (kattarikkāy) — an aubergine, a brinjal}

\label{lesson:TA-C243-kattarikkay}



\begin{quote}
Before the new one: say the Tamil for wheat, then the Tamil for a potato.
\end{quote}

\subsection*{You'll want to know: \ta{கத்தரிக்காய்}}
\textbf{\ta{கத்தரிக்காய்}} — \emph{kattarikkāy} — \textquotedblleft{}an aubergine, a brinjal\textquotedblright{}.

\begin{grammarlens}[title={What you've built}]
One more word for this chapter.
\end{grammarlens}

\subsection*{Your turn}
\begin{itemize}
\raggedright
  \item \emph{Say it:} \emph{kattarikkāy}
  \item \emph{Say it:} \emph{kattarikkāy}, once more
  \item \emph{Say it:} say \emph{kattarikkāy}
  \item \emph{Recall:} say the Tamil for wheat, then the Tamil for a potato, then say \emph{kattarikkāy} again
\end{itemize}

\begin{tcolorbox}[breakable,colback=teal!4,colframe=teal!35!black,title={Before you move on}]
What does \ta{கத்தரிக்காய்} mean? (\textquotedblleft{}An aubergine, a brinjal\textquotedblright{}.) Say it once more.
\end{tcolorbox}
\section[\texorpdfstring{\ta{வாசனை} (vācaṉai)}{வாசனை (vācaṉai)}]{\ta{வாசனை} (vācaṉai) — a fragrance, a smell}

\label{lesson:TA-C243-vacanai}



\begin{quote}
Before the new one: say the Tamil for a potato, then the Tamil for an aubergine.
\end{quote}

\subsection*{You'll want to know: \ta{வாசனை}}
\textbf{\ta{வாசனை}} — \emph{vācaṉai} — \textquotedblleft{}a fragrance, a smell\textquotedblright{}.

\begin{grammarlens}[title={What you've built}]
One more word for this chapter.
\end{grammarlens}

\subsection*{Your turn}
\begin{itemize}
\raggedright
  \item \emph{Say it:} \emph{vācaṉai}
  \item \emph{Say it:} \emph{vācaṉai}, once more
  \item \emph{Say it:} say \emph{vācaṉai}
  \item \emph{Recall:} say the Tamil for a potato, then the Tamil for an aubergine, then say \emph{vācaṉai} again
\end{itemize}

\begin{tcolorbox}[breakable,colback=teal!4,colframe=teal!35!black,title={Before you move on}]
What does \ta{வாசனை} mean? (\textquotedblleft{}A fragrance, a smell\textquotedblright{}.) Say it once more.
\end{tcolorbox}
\section[\texorpdfstring{\ta{வெள்ளரிக்காய்} (veḷḷarikkāy)}{வெள்ளரிக்காய் (veḷḷarikkāy)}]{\ta{வெள்ளரிக்காய்} (veḷḷarikkāy) — a cucumber}

\label{lesson:TA-C243-vellarikkay}



\begin{quote}
Before the new one: say the Tamil for an aubergine, then the Tamil for a fragrance.
\end{quote}

\subsection*{You'll want to know: \ta{வெள்ளரிக்காய்}}
\textbf{\ta{வெள்ளரிக்காய்}} — \emph{veḷḷarikkāy} — \textquotedblleft{}a cucumber\textquotedblright{}.

\begin{grammarlens}[title={What you've built}]
One more word for this chapter.
\end{grammarlens}

\subsection*{Your turn}
\begin{itemize}
\raggedright
  \item \emph{Say it:} \emph{veḷḷarikkāy}
  \item \emph{Say it:} \emph{veḷḷarikkāy}, once more
  \item \emph{Say it:} say \emph{veḷḷarikkāy}
  \item \emph{Recall:} say the Tamil for an aubergine, then the Tamil for a fragrance, then say \emph{veḷḷarikkāy} again
\end{itemize}

\begin{tcolorbox}[breakable,colback=teal!4,colframe=teal!35!black,title={Before you move on}]
What does \ta{வெள்ளரிக்காய்} mean? (\textquotedblleft{}A cucumber\textquotedblright{}.) Say it once more.
\end{tcolorbox}
